\documentclass{article}
\usepackage[UTF8]{ctex}
\usepackage{amsmath,amssymb,amsfonts}
\usepackage{enumitem}
\usepackage{geometry}
\geometry{a4paper, margin=1in}
% === 错题本排版优化（不参与题目标识） ===
\setlength{\parindent}{0pt}
\setlength{\parskip}{0.35em plus 0.10em minus 0.08em}
\setlength{\emergencystretch}{3em}
\setlist{itemsep=0.25em, topsep=0.40em, parsep=0pt}
\allowdisplaybreaks[2]
% === 错题本排版优化结束 ===
\begin{document}
\section*{东北师范大学2026年高等代数真题}
\section*{73．东北师范大学2026年高等代数与解析几何试卷}

\subsection{题目1}
1．（20 分）矩阵 \(A=\left(\begin{array}{cccc}1 & 1 & 1 & 1 \\ 1 & -1 & -1 & 1 \\ 1 & -1 & 1 & -1 \\ 1 & 1 & -1 & -1\end{array}\right)\) ，求 \(A^{2026}\) ．

\subsection{题目2}
2．（10分）求同余方程组

\[
\begin{cases}x \equiv 1 & (\bmod 4) \\ x \equiv-3 & (\bmod 7) \\ x \equiv 2 & (\bmod 15)\end{cases}
\]

\subsection{题目3}
3．（15 分）已知 \(A\) 是 \(m \times n\) 阶矩阵，\(B\) 是 \(n \times m\) 阶矩阵。证明：

\[
\left|\begin{array}{cc}
E_{n} & B \\
A & E_{m}
\end{array}\right|=\left|E_{m}-A B\right|=\left|E_{n}-B A\right| .
\]

\subsection{题目4}
4．（15 分）已知 \(\varepsilon_{1}, \varepsilon_{2}, \cdots, \varepsilon_{n}\) 是 \(n\) 维欧氏空间 \(V\) 的一组标准正交基，对任意的 \(\alpha, \beta \in V\) ，证明：

\[
(\alpha, \beta)=\sum_{i=1}^{n}\left(\alpha, \varepsilon_{i}\right)\left(\beta, \varepsilon_{i}\right) .
\]

\subsection{题目5}
5．（20 分）对于有限维线性空间 \(V, \mathcal{A}\) 是 \(V\) 上的一个线性变换，证明：存在 \(V\) 上的两个线性变换 \(\mathcal{D}, \mathcal{N}\)，满足

\begin{itemize}
  \item[（1）] 9 可对角化．
  \item[（2）] \(\mathcal{N}\) 为幂零变换．
  \item[（3）] \(\mathcal{A}=\mathcal{D}+\mathcal{N}, \mathcal{D} \mathcal{N}=\mathcal{N} \mathcal{D}\) ．
\end{itemize}

\subsection{题目6}
6．（20 分）设 \(\varphi, \theta\) 是线性空间 \(V\) 上的两个线性变换，且 \(f(x), g(x)\) 分别为 \(\varphi, \theta\) 的特征多项式。

\begin{itemize}
  \item[（1）] 若 \(h(x)\) 是一个多项式，\(a\) 为 \(\varphi\) 的一个特征值，证明：\(h(a)\) 是 \(h(\varphi)\) 的一个特征值．
  \item[（2）] 若 \(f(x)\) 与 \(g(x)\) 互素，证明：\(f(\theta)\) 和 \(g(\varphi)\) 都是可逆的线性变换．
\end{itemize}

\subsection{题目7}
7．（15 分）已知空间曲线 \(L\) 的参数方程为 \(\left\{\begin{array}{l}x=f(t) \\ y=g(t) \\ z=h(t)\end{array}\right.\) ，在空间直角坐标系下，将曲线 \(L\) 绕 \(y\) 轴旋转，求旋转曲面方程。

\subsection{题目8}
8．（15 分）已知直线 \(\left\{\begin{array}{l}x-y+z-1=0, \\ 2 x-y+z-2=0 .\end{array}\right.\) 平面 \(\Pi: x+y+z-2=0\) ．

\begin{itemize}
  \item[（1）] 求直线 \(L\) 与平面 \(\Pi\) 夹角的正弦值．
  \item[（2）] 求直线 \(L\) 在平面 \(\Pi\) 上的射影直线方程。
\end{itemize}

\subsection{题目9}
9．（20 分）已知直线

\[
L_{1}:\left\{\begin{array}{l}
x-2 y-2=0, \\
y-z+2=0 .
\end{array} L_{2}:\left\{\begin{array}{l}
x-2 y-2=0, \\
x-z+1=0 .
\end{array}\right.\right.
\]

平面 \(\Pi: x-y+2=0\) ．求与直线 \(L_{1}, L_{2}\) 都相交且与平面 \(\Pi\) 垂直的直线的轨迹方程．
\end{document}
